\section*{ЗАКЛЮЧЕНИЕ}
\addcontentsline{toc}{section}{ЗАКЛЮЧЕНИЕ}
В процессе выполнения работы была создана программно-информационная система «Цифровая платформа Курский край» — интерактивный веб-гид по городу Курску. Были реализованы функции для пользователей и администратора.

Пользователи имеют возможность просматривать список мест с фильтрацией по категориям (парки, музеи, кафе, отели и т.д.), открывать подробные карточки объектов с описанием, контактами и галереей, а также просматривать маршруты и видеть их визуализацию на карте. Администратор может добавлять, редактировать и удалять объекты и маршруты, а также создавать новые категории.

Основные результаты работы:

\begin{enumerate}
	\item Проведён анализ предметной области. Выявлены ключевые сущности (категории, места, маршруты, точки маршрута) и обоснована необходимость автоматизации.
	\item Разработана концептуальная модель веб-сайта. Построена ER-модель базы данных. Определены функциональные и нефункциональные требования к системе.
	\item Осуществлено проектирование веб-сайта. Разработана архитектура клиент-сервер. Спроектирован пользовательский интерфейс.
	\item Реализован и протестирован веб-сайт. Проведено модульное и системное тестирование.
\end{enumerate}

Все требования, объявленные в техническом задании, были полностью реализованы, все задачи, поставленные в начале разработки проекта, также решены.

Сайт находится в публичном доступе, поскольку опубликован в сети Интернет.  
